% vim: ai sw=2 spell spelllang=cs

\documentclass[xcolor=dvipsnames,handout]{beamer}


\usepackage[czech]{babel}
\uselanguage{czech}
\deftranslation[to=czech]{Section}{Sekce}

\let\Oldemph\emph
\renewcommand{\emph}[1]{\textcolor{blue}{\Oldemph{#1}}}
\newcommand{\heading}[1]{\textbf{\textcolor{BurntOrange}{#1}}}
\newcommand{\header}[1]{\textbf{\texttt{\textcolor{WildStrawberry}{#1}}}}
\newcommand{\doc}[1]{\textbf{\texttt{\textcolor{OliveGreen}{#1}}}}
\newcommand{\red}[1]{\textcolor{red}{#1}}

\mode<presentation>
{
\usetheme{Madrid}
}

\setbeamertemplate{navigation symbols}{}

\usepackage{adjustbox}
\usepackage[utf8x]{inputenc}
\usepackage{helvet}
\usepackage{mathptmx}
\usepackage[T1]{fontenc}
\usepackage{graphicx}
\usepackage{hyperref}
\usepackage[final]{listings}
\usepackage{multirow}
\usepackage{tikz}
\usepackage[absolute,overlay]{textpos}
\usepackage{xspace}
\usetikzlibrary{arrows, positioning}

\tikzset{tl_state/.style={rounded corners, thick, draw, minimum height=6mm,
	minimum width=12mm, bottom color=black!15, top color=black!5}}
\tikzset{tl_phase/.style={thick, draw, minimum size=8mm}}
\tikzset{tl_edge/.style={<-, >=stealth', semithick}}

\lstset{tabsize=2,numbers=none,showstringspaces=false,breaklines=true,
	breakatwhitespace=true,escapeinside={(*@}{@*)},
	basicstyle=\scriptsize,keywordstyle=\bfseries\color{OliveGreen},
	commentstyle=\color{blue}, columns=flexible, language=C,
	morekeywords={asm}}

\newcommand{\cmd}[2][]{\lstinline[columns=fullflexible, language=bash,
	basicstyle=\ttfamily, #1]$#2$}
\newcommand{\code}[2][]{\lstinline[columns=fullflexible, basicstyle=\ttfamily,
	#1]$#2$}


\newcommand{\make}{\texttt{make}\xspace}
\newcommand{\cmake}{\texttt{cmake}\xspace}

\title[PB071 -- Překlad]{PB071 -- Úvod do jazyka C}

\author{Jiri~Slaby}

\institute[Fakulta informatiky, MU]{Fakulta informatiky \\
Masarykova univerzita}

\subtitle{Překladové systémy}
\date{24. 4. \the\year}

\begin{document}

\begin{frame}
  \titlepage
\end{frame}

\begin{frame}{Obsah přednášky}
  \tableofcontents
\end{frame}

%%%%%%%%%%%%%%%%%%%%%%%

\section{Úvod a motivace}
\frame{\sectionpage}

\begin{frame}{Úvod}

  \begin{itemize}
    \item Máme zdrojové soubory

    \pause

    \item Chceme binárku

    \pause

    \item Použijeme překladač
    \begin{itemize}
      \item \cmd{gcc *.c}
    \end{itemize}

    \pause

    \item \heading{Demo:} jádro

    \pause

    \item Chceme raději nějaký nástroj
    \begin{itemize}
      \item Popis, jak sestavit projekt

      \pause

      \item Ale jaký?
    \end{itemize}

  \end{itemize}

\end{frame}

\begin{frame}{Překladové systémy}

  \heading{Spousta řešení}
  \begin{itemize}

    \item Svoje vlastní
    \begin{itemize}
      \item Skripty v bashi, perlu, pythonu

      \pause

      \item \emph{Vždy} peklo

      \pause

      \begin{itemize}
	\item \emph{Nikdy} nevynalézejte kolo
      \end{itemize}
    \end{itemize}

    \pause

    \item Cizí
    \begin{itemize}
      \item Udržované nástroje

      \pause

      \item \make, \cmake, \texttt{qmake}, \dots
    \end{itemize}

  \end{itemize}

\end{frame}

\begin{frame}{Obecný princip}

\begin{center}
  \heading{Popis závislostí + popis výroby} \\[1ex]

  \begin{textblock*}{3.2cm}(\textwidth-3.2cm,1.2cm)
    \begin{onlyenv}<3->
      \tikzstyle{pre} = [<-,shorten <=1pt,>=stealth',semithick,above]
      \begin{tikzpicture}[node distance=7mm, font=\footnotesize]
	\node[tl_state] (X) {X};
	\node[tl_state] (Y) [right=of X] {Y}
	edge[pre] node (MID) {} (X);
	\node (DESC) [below=5mm of MID] {X závisí na Y};
      \end{tikzpicture}
    \end{onlyenv}
  \end{textblock*}

  \tikzstyle{pre} = [<-,shorten <=1pt,>=stealth',semithick,above]
  \tikzstyle{post} = [pre,->,shorten >=1pt, shorten <=0pt]
  \begin{tikzpicture}[node distance=7mm, font=\footnotesize]
      \onslide<6->{
	\node [tl_state] (BIN) {\file{binary}};
      }

      \onslide<5->{
	\node [tl_state] (O2) [below=of BIN] {\file{impl.o}};
	\node [tl_state] (O1) [left=of O2] {\file{main.o}};
      }

      \onslide<7->{
	\node [tl_state] (LIB1) [right=of O2] {\file{libX.so}}
	  edge[pre] (BIN);
      }

      \onslide<6->{
	\draw[pre] (O2) -- (BIN);
	\draw[pre] (O1) -- node (LDP) {} (BIN);
      }

      \node [tl_state] (C2) [below=of O2] {\file{impl.c}};
      \node [tl_state] (C1) [below=of O1] {\file{main.c}};

      \onslide<5-> {
	\draw[pre] (C2) -- (O2);
	\draw[pre] (C1) -- node[yshift=-1mm] (GCCP) {} (O1) ;
      }

      \onslide<2->{
	\node [tl_state] (H1) [below=of C1] {\file{stdio.h}};
	\node [tl_state] (H2) [below=of C2] {\file{impl.h}};
      }

      \onslide<3->{
	\draw[pre] (H1) -- (C1);
	\draw[pre] (H2) -- (C2);
	\draw[pre] (H2) -- (C1);
      }

      \onslide<4->{
	\draw [pre,out=310,in=230] (H1) edge (H2);
	\node [tl_state] (SW) [below=of H2,align=center]
	  {Parametry překladače \\
	  (\cmd{\-D apod.})}
	  edge [pre,dashed,out=40,in=340] node {} (C2)
	  edge [pre,dashed] node {} (H2);
      }

      \onslide<8->{
        \node (GCC) [left=of GCCP] {\cmd{gcc}}
	  edge [post,dotted,color=red] node {} (GCCP);

        \node (LD) [left=of LDP] {\cmd{ld}}
	  edge [post,dotted,color=red] node {} (LDP);

      }

      \onslide<9->{
        \node (EX) [below=5mm of LIB1,xshift=8mm,align=center] {Musí \\ existovat}
	  edge [post,dotted,color=red,bend right=25] node {} (C1)
	  edge [post,dotted,color=red,bend right=10] node {} (C2)
	  edge [post,dotted,color=red,bend right=5] node {} (H1)
	  edge [post,dotted,color=red,bend left=15] node {} (H2)
	  edge [post,dotted,color=red] node {} (LIB1);
      }

  \end{tikzpicture}
\end{center}

\end{frame}

\section{Vývojový nástroj make}
\frame{\sectionpage}

\begin{frame}{Vývojový nástroj \make}

  \begin{itemize}
    \item Nejrozšířenější v UN*Xovém světě

    \pause

    \item Pravidla pro překlad projektu
    \begin{itemize}
      \item V textové formě
    \end{itemize}
    \item \cmd{make} pravidla načte, vyhodnotí a spustí

    \pause

    \item Mezi pravidly jsou typicky závislosti
    \begin{itemize}
      \item Viz předchozí slajd
    \end{itemize}

    \pause

    \item \cmd{make} hledá \file{Makefile}, popř.\ \file{makefile} v lokálním
      adresáři
    \begin{itemize}
      \item \cmd{\-f} mění název souboru
      \item \cmd{\-C} mění adresář
    \end{itemize}

    \pause

    \item \url{https://www.gnu.org/software/make/manual/make.html}
  \end{itemize}

\end{frame}

\begin{frame}[fragile]{Makefile}{Pravidla}

  \begin{itemize}
    \item Pravidla: \code{cil: zavislosti}
    \begin{itemize}
      \item Implicitně se začne vyhodnocením prvního v souboru

      \pause

      \item Ale lze specifikovat na příkazové řádce: \cmd{make cil}
    \end{itemize}

    \pause

    \item Pravidla můžou mít akce, jak \code{cil} vyrobit ze
      \code{zavislosti}
    \begin{itemize}
      \item Na dalším řádku za pravidlem, \emph{odsazené tabulátorem}

      \pause

      \item Volá se shell

      \pause

      \item Jinak se použije implicitní akce
    \end{itemize}
  \end{itemize}

  \begin{onslide}<1->
  \begin{block}{Příklad}
  \begin{lstlisting}[language=make]
all: binary

binary: main.o impl.o
       gcc -obinary main.o impl.o
  \end{lstlisting}
  \end{block}
  \end{onslide}

\end{frame}

\begin{frame}{Úkol}

  \heading{Hello world Makefile}
  \begin{enumerate}
    \item Napište jednoduchý Makefile
    \item V něm budou 3 pravidla
    \begin{itemize}
      \item \code{p1}, \code{p2}, \code{p3}
    \end{itemize}
    \item Každé vypíše své jméno
    \item Každé z nich vyvolejte
    \begin{itemize}
      \item Z příkazové řádky
    \end{itemize}

    \item Nyní mezi nimi vytvořte závislosti
    \begin{itemize}
      \item Ať se vypíše po řadě \code{p1}, \code{p2}, \code{p3}
    \end{itemize}
    \item Nyní spusťte jen samotný \cmd{make}
  \end{enumerate}

\end{frame}

\begin{frame}[fragile]{Makefile}{Proměnné}

  \begin{itemize}
    \item Nastavení proměnných
    \begin{itemize}
      \item Na příkazové řádce
      \item V souboru: \code{CC=gcc}
    \end{itemize}

    \pause

    \item Reference pomocí \code{\$(...)}
  \end{itemize}

  \begin{block}{Příklad}
  \begin{lstlisting}[language=make]
CFLAGS=-Wall

all: binary

binary: main.o impl.o
       $(CC) -obinary main.o impl.o
  \end{lstlisting} %stopzone
  \end{block}

\end{frame}

\begin{frame}{Úkol}

  \heading{Proměnné v Makefile}
  \begin{enumerate}
    \item Vytvořte nové pravidlo \file{m.o}
    \item Bude záviset na \file{main.c}
    \item Bude překládat \file{main.c} na \file{m.o}
    \item Bude brát v úvahu proměnné \code{CC} a \code{CFLAGS}
    \item Upravte jedno z předchozích \code{p*} pravidel, aby záviselo na
      \file{m.o}
    \item Vyzkoušejte
    \begin{itemize}
      \item Zkuste předat různé \code{CFLAGS}
      \item Nezapomeňte pozměnit alespoň čas souboru (\cmd{touch main.c})
    \end{itemize}
  \end{enumerate}

\end{frame}

\begin{frame}[fragile]{Makefile}{Obecná pravidla}

  \begin{itemize}
    \item \code{\%} je zástupný symbol
    \item Lze psát pravidla: \code{\%.o: \%.c}

    \pause

    \item Jak ale zjistit jména souborů?
    \begin{itemize}
      \item \code{\$@} -- jméno cíle
      \item \code{\$<} -- jméno první závislosti
      \item \code{\$^} -- jména všech závislostí
    \end{itemize}
  \end{itemize}

  \pause

  \begin{block}{Příklad}
  \begin{lstlisting}[language=make]
%.o: %.c
       $(CC) -c -o$@ $<
  \end{lstlisting} %stopzone
  \end{block}

\end{frame}

\begin{frame}{Úkol}

  \heading{Obecná pravidla v Makefile}
  \begin{enumerate}
    \item Vytvořte nové pravidlo k překladu z \file{.c} do \file{.o}
    \item Bude brát v úvahu proměnné \code{CC} a \code{CFLAGS}
    \item Kromě \code{CFLAGS} bude překládat s \cmd{\-W}
    \item Upravte jedno z předchozích \code{p*} pravidel, aby záviselo na
      \file{main.o}
    \item Vyzkoušejte
  \end{enumerate}

\end{frame}

\begin{frame}[fragile]{Makefile}{Vnitřní příkazy}

  \begin{itemize}
    \item Make má vnitřní příkazy
    \begin{itemize}
      \item \cmd{\$(prikaz parametry)}
    \end{itemize}

    \pause

    \item \code{shell command}
    \begin{itemize}
      \item Spustí shell a vykoná v něm příkaz \code{command}
    \end{itemize}

    \pause

    \item \cmd{wildcard pattern}
    \begin{itemize}
      \item Expanduje \cmd{pattern} (\cmd{*} ne \cmd{\%}) a vrací seznam
    \end{itemize}

    \pause

    \item \cmd{patsubst pattern,replacement,text}
    \begin{itemize}
      \item Nahraď v \cmd{text} všechny výskyty \cmd{pattern} za
        \cmd{replacement}
    \end{itemize}

    \pause

    \item \dots
  \end{itemize}

  \onslide<2->

  \begin{block}{Příklad}
    \begin{lstlisting}[language=make,belowskip=0pt]
PATH=$(shell pwd)
    \end{lstlisting} %stopzone
    \onslide<3->
    \begin{lstlisting}[language=make,aboveskip=0pt,belowskip=0pt]
SOURCES=$(wildcard *.c)
    \end{lstlisting} %stopzone
    \onslide<4->
    \begin{lstlisting}[language=make,aboveskip=0pt]
OBJS=$(patsubst %.c,%.o,$(SOURCES))
    \end{lstlisting} %stopzone
  \end{block}

\end{frame}

\begin{frame}{Makefile}{Nevýhody}

  \begin{itemize}
    \item Nutno popsat všechny závislosti

    \pause

    \item Podpora pro \cmd{make clean} ručně

    \pause

    \item Rychlost
    \begin{itemize}
      \item Není lineární vzhledem k počtu souborů
    \end{itemize}

    \pause

    \item Složitý zápis komplikovanějších pravidel
    \begin{itemize}
      \item Např.\ neexistuje explicitní podpora pro podadresáře
    \end{itemize}

    \pause

    \item Méně přenositelný

    \pause

    \item Může být příliš upovídaný
  \end{itemize}

\end{frame}

\section{Autotools, CMake}
\frame{\sectionpage}

\begin{frame}{Generátory}

  \begin{itemize}
    \item Nadstavby nad \make (a ostatními)
    \begin{itemize}
      \item Často generují \file{Makefile} pro \make
    \end{itemize}

    \pause

    \item Podpora externích závislostí
    \begin{itemize}
      \item Je v dosahu překladač C?
      \item Umí překladač C99?
      \item Je v systému knihovna X?
      \item Máme hlavičku \file{Y.h}?
    \end{itemize}

    \pause

    \item Podpora podadresářů

    \pause

    \item Konfigurace uživatelem
    \begin{itemize}
      \item Ne/chci ve výsledné binárce podporu pro X
      \item Ne/optimalizuj při překladu (ne/chci ladit)
    \end{itemize}

    \pause

    \item \dots
  \end{itemize}

\end{frame}

\begin{frame}{Nástroje Autotools}

  \heading{Tři nástroje}
  \begin{enumerate}
    \item \texttt{Autoconf}
    \begin{itemize}
      \item Vyhledání externích závislostí
      \item Konfigurace uživatelem
      \item Z popisu v \file{configure.ac} generuje skript \file{configure}
      \item \url{http://www.gnu.org/software/autoconf/manual/}
    \end{itemize}

    \pause

    \item \texttt{Automake}
    \begin{itemize}
      \item Popis překladu (včetně rozdělení do podadresářů)
      \item Z popisu v \file{Makefile.am} generuje \file{Makefile.in}
      \item \url{http://www.gnu.org/software/automake/manual/}
    \end{itemize}

    \pause

    \item \texttt{Libtool}
    \begin{itemize}
      \item Podpora pro tvorbu knihoven na různých POSIX systémech
      \item \url{https://www.gnu.org/software/libtool/manual/}
    \end{itemize}
  \end{enumerate}

  \pause
  \medskip

  \heading{Výsledek: jeden \file{configure} a mnoho \file{Makefile.in}}

  \pause
  \smallskip

  \heading{Uživatel spustí \cmd{./configure} a dostane mnoho \file{Makefile}}

\end{frame}

\begin{frame}[fragile]{Překladový systém \cmake}

  \begin{itemize}
    \item Popis konfigurace a překladu v jednom
    \begin{itemize}
      \item Soubory \file{CMakeLists.txt}
      \item Znáte z cvičení (\file{kontr_lessons})
    \end{itemize}

    \pause

    \item \url{https://cmake.org/cmake/help/latest/}
  \end{itemize}

  \pause

  \begin{block}{Příklad}
  \begin{lstlisting}[language=make,columns=fullflexible]
cmake_minimum_required(VERSION 2.8.11)
project(my_project)

set(CMAKE_C_FLAGS "-std=c99 -pedantic -Wall -Wextra")

add_subdirectory(my_lib)
add_executable(main main.c impl.c)
target_link_libraries(main LINK_PUBLIC my_lib)
  \end{lstlisting}
  \end{block}

\end{frame}

\section{Závěr}
\frame{\sectionpage}

\begin{frame}{Závěr a diskuse}

  \begin{center} \Large
    \heading{Nepište překladové systémy na vlastní pěst} \\[2em]

    \pause \huge

    Děkuji za pozornost! \\[1em]
    Dotazy?
  \end{center}

\end{frame}

\end{document}
